%
% $Id: slides.tex 4228 2006-06-21 21:55:12Z jjamor $
%
%
% Compilar a .pdf con LaTeX (pdflatex)
% Es necesario instalar Beamer (paquete latex-beamer en Debian)
%

%
% Gr�ficos:
% Los gr�ficos pueden suministrarse en PNG, JPG, TIF, PDF, MPS
% Los EPS deben convertirse a PDF (usar epstopdf)
%

\documentclass{beamer}
\usetheme{Warsaw}
\usebackgroundtemplate{\includegraphics[width=\paperwidth]{format/libresoft-bg.png}}
\usepackage[spanish]{babel}
\usepackage[latin1]{inputenc}
\usepackage{graphics}
\usepackage{amssymb} % Simbolos matematicos
\usepackage{url}

%\definecolor{libresoftgreen}{RGB}{162,190,43}
%\definecolor{libresoftblue}{RGB}{0,98,143}

%\setbeamercolor{titlelike}{bg=libresoftgreen}

%% Metadatos del PDF.
\hypersetup{
  pdftitle={Tools for Information and Documentation Management},
  pdfauthor={Daniel Izquierdo Cort�zar},
  pdfcreator={GSyC/Libresoft},
  pdfproducer=PDFLaTeX,
  pdfsubject={},
}
%%


\AtBeginSection[]
{
  \begin{frame}<presentation>
    \frametitle{Index}
    \tableofcontents[current]
  \end{frame}
}


\begin{document}

\title{Tools for Information and Documentation Management}
\subtitle{Master on Free Software}
\institute{dizquierdo@gsyc.escet.urjc.es\\
GSyC/Libresoft}
\author{Daniel Izquierdo Cort�zar}
\date{\today}

\frame{
\maketitle
\begin{center}
\includegraphics[width=6cm]{format/gsyc-urjc}
\end{center}
}


% Si el titulo o el autor se quieren acortar para los pies de p�gina
% se pueden redefinir aqu�:
%\title{Titulo corto}
%\author{Autores abreviado}


%% LICENCIA DE REDISTRIBUCION DE LAS TRANSPAS
\frame{
~
\vspace{4cm}

\begin{flushright}
{\tiny
(cc) 2007 Daniel Izquierdo Cort�zar. \\
Some rights reserved. This document is distributed under the Creative \\
            Commons Attribution-ShareAlike 2.5 licence, available in \\
            http://creativecommons.org/licenses/by-sa/2.5/

%  Este documento (o uno muy similar) está disponible en \\
%  \url{http://gsyc.escet.urjc.es/~jjamor/}
}
\end{flushright}
}
%%



\section{Introduction}

\begin{frame}
\frametitle{}
\begin{itemize}
\item Based on QualipSO project (\url{http://www.qualipso.org}) \\
 
\begin{quotation}
``QualiPSo is a unique alliance of European, Brazilian and Chinese ICT 
industry players, SMEs, governments and academics to help industries 
and governments fuel innovation and competitiveness with Open Source 
software.''
\end{quotation}
\item Deliverable made by Libresoft (\url{http://libresoft.es})
\end{itemize}
\end{frame}

\section{Communication and Coordination}

\begin{frame}
\frametitle{Tools}
\begin{itemize}
\item Mailman
\item SmartList
\item XChat
\item Gaim
\item Ekiga
\item MoinMoin
\item Evolution
\item Planner
\item Trac
\item Others
\end{itemize}
\end{frame}

\section{Documentation}
\begin{frame}
\frametitle{Tools}
\begin{itemize}
\item Plone
\item Joomla!
\item Open Office
\item KOffice
\item Latex
\item Dia
\item FreeMind
\item Others
\end{itemize}
\end{frame}



\section{Development}
\begin{frame}
\frametitle{Tools}
\begin{itemize}
\item CVS, SVN, Git, Bazaar, ...
\item Emacs, VI, ...
\item Bugzilla, debbugs
\item Glade, Stetic, ...
\item apt, rpm, yum, ...
\end{itemize}
\end{frame}

\section{Support}
\begin{frame}
\frametitle{Tools}
\begin{itemize}
\item Bugzilla
\item phpBB
\item Mailman
\item XChat
\item Others
\end{itemize}
\end{frame}

\section{Translation}
\begin{frame}
\frametitle{Tools}
\begin{itemize}
\item gettext
\item KBabel
\item gtranslator
\item Open Language Tools
\item xml2po
\item Others
\end{itemize}
\end{frame}


\section{Bibiography}
\begin{frame}
\begin{itemize}

\item �lvaro del Castillo, Julia Anaya. QualipSO D4.2.2: 
Available tools for Information and Documentation Management.
Technical Report, Universidad Rey Juan Carlos, M�stoles, Spain, May 2007.
\url{http://www.qualipso.org/media/A4/D4.2.2.pdf}. Last visited November, 2007.

\item Karl Fogel. Producing Open Source Software: How to run a succesful
free software project. Chapter 8: Managing Volunteers. Ed: O'Reilly, October 2005.

\end{itemize}
\end{frame}





\end{document}
